% Options for packages loaded elsewhere
\PassOptionsToPackage{unicode}{hyperref}
\PassOptionsToPackage{hyphens}{url}
%
\documentclass[
]{article}
\usepackage{amsmath,amssymb}
\usepackage{iftex}
\ifPDFTeX
  \usepackage[T1]{fontenc}
  \usepackage[utf8]{inputenc}
  \usepackage{textcomp} % provide euro and other symbols
\else % if luatex or xetex
  \usepackage{unicode-math} % this also loads fontspec
  \defaultfontfeatures{Scale=MatchLowercase}
  \defaultfontfeatures[\rmfamily]{Ligatures=TeX,Scale=1}
\fi
\usepackage{lmodern}
\ifPDFTeX\else
  % xetex/luatex font selection
\fi
% Use upquote if available, for straight quotes in verbatim environments
\IfFileExists{upquote.sty}{\usepackage{upquote}}{}
\IfFileExists{microtype.sty}{% use microtype if available
  \usepackage[]{microtype}
  \UseMicrotypeSet[protrusion]{basicmath} % disable protrusion for tt fonts
}{}
\makeatletter
\@ifundefined{KOMAClassName}{% if non-KOMA class
  \IfFileExists{parskip.sty}{%
    \usepackage{parskip}
  }{% else
    \setlength{\parindent}{0pt}
    \setlength{\parskip}{6pt plus 2pt minus 1pt}}
}{% if KOMA class
  \KOMAoptions{parskip=half}}
\makeatother
\usepackage{xcolor}
\usepackage[margin=1in]{geometry}
\usepackage{color}
\usepackage{fancyvrb}
\newcommand{\VerbBar}{|}
\newcommand{\VERB}{\Verb[commandchars=\\\{\}]}
\DefineVerbatimEnvironment{Highlighting}{Verbatim}{commandchars=\\\{\}}
% Add ',fontsize=\small' for more characters per line
\usepackage{framed}
\definecolor{shadecolor}{RGB}{248,248,248}
\newenvironment{Shaded}{\begin{snugshade}}{\end{snugshade}}
\newcommand{\AlertTok}[1]{\textcolor[rgb]{0.94,0.16,0.16}{#1}}
\newcommand{\AnnotationTok}[1]{\textcolor[rgb]{0.56,0.35,0.01}{\textbf{\textit{#1}}}}
\newcommand{\AttributeTok}[1]{\textcolor[rgb]{0.13,0.29,0.53}{#1}}
\newcommand{\BaseNTok}[1]{\textcolor[rgb]{0.00,0.00,0.81}{#1}}
\newcommand{\BuiltInTok}[1]{#1}
\newcommand{\CharTok}[1]{\textcolor[rgb]{0.31,0.60,0.02}{#1}}
\newcommand{\CommentTok}[1]{\textcolor[rgb]{0.56,0.35,0.01}{\textit{#1}}}
\newcommand{\CommentVarTok}[1]{\textcolor[rgb]{0.56,0.35,0.01}{\textbf{\textit{#1}}}}
\newcommand{\ConstantTok}[1]{\textcolor[rgb]{0.56,0.35,0.01}{#1}}
\newcommand{\ControlFlowTok}[1]{\textcolor[rgb]{0.13,0.29,0.53}{\textbf{#1}}}
\newcommand{\DataTypeTok}[1]{\textcolor[rgb]{0.13,0.29,0.53}{#1}}
\newcommand{\DecValTok}[1]{\textcolor[rgb]{0.00,0.00,0.81}{#1}}
\newcommand{\DocumentationTok}[1]{\textcolor[rgb]{0.56,0.35,0.01}{\textbf{\textit{#1}}}}
\newcommand{\ErrorTok}[1]{\textcolor[rgb]{0.64,0.00,0.00}{\textbf{#1}}}
\newcommand{\ExtensionTok}[1]{#1}
\newcommand{\FloatTok}[1]{\textcolor[rgb]{0.00,0.00,0.81}{#1}}
\newcommand{\FunctionTok}[1]{\textcolor[rgb]{0.13,0.29,0.53}{\textbf{#1}}}
\newcommand{\ImportTok}[1]{#1}
\newcommand{\InformationTok}[1]{\textcolor[rgb]{0.56,0.35,0.01}{\textbf{\textit{#1}}}}
\newcommand{\KeywordTok}[1]{\textcolor[rgb]{0.13,0.29,0.53}{\textbf{#1}}}
\newcommand{\NormalTok}[1]{#1}
\newcommand{\OperatorTok}[1]{\textcolor[rgb]{0.81,0.36,0.00}{\textbf{#1}}}
\newcommand{\OtherTok}[1]{\textcolor[rgb]{0.56,0.35,0.01}{#1}}
\newcommand{\PreprocessorTok}[1]{\textcolor[rgb]{0.56,0.35,0.01}{\textit{#1}}}
\newcommand{\RegionMarkerTok}[1]{#1}
\newcommand{\SpecialCharTok}[1]{\textcolor[rgb]{0.81,0.36,0.00}{\textbf{#1}}}
\newcommand{\SpecialStringTok}[1]{\textcolor[rgb]{0.31,0.60,0.02}{#1}}
\newcommand{\StringTok}[1]{\textcolor[rgb]{0.31,0.60,0.02}{#1}}
\newcommand{\VariableTok}[1]{\textcolor[rgb]{0.00,0.00,0.00}{#1}}
\newcommand{\VerbatimStringTok}[1]{\textcolor[rgb]{0.31,0.60,0.02}{#1}}
\newcommand{\WarningTok}[1]{\textcolor[rgb]{0.56,0.35,0.01}{\textbf{\textit{#1}}}}
\usepackage{longtable,booktabs,array}
\usepackage{calc} % for calculating minipage widths
% Correct order of tables after \paragraph or \subparagraph
\usepackage{etoolbox}
\makeatletter
\patchcmd\longtable{\par}{\if@noskipsec\mbox{}\fi\par}{}{}
\makeatother
% Allow footnotes in longtable head/foot
\IfFileExists{footnotehyper.sty}{\usepackage{footnotehyper}}{\usepackage{footnote}}
\makesavenoteenv{longtable}
\usepackage{graphicx}
\makeatletter
\def\maxwidth{\ifdim\Gin@nat@width>\linewidth\linewidth\else\Gin@nat@width\fi}
\def\maxheight{\ifdim\Gin@nat@height>\textheight\textheight\else\Gin@nat@height\fi}
\makeatother
% Scale images if necessary, so that they will not overflow the page
% margins by default, and it is still possible to overwrite the defaults
% using explicit options in \includegraphics[width, height, ...]{}
\setkeys{Gin}{width=\maxwidth,height=\maxheight,keepaspectratio}
% Set default figure placement to htbp
\makeatletter
\def\fps@figure{htbp}
\makeatother
\setlength{\emergencystretch}{3em} % prevent overfull lines
\providecommand{\tightlist}{%
  \setlength{\itemsep}{0pt}\setlength{\parskip}{0pt}}
\setcounter{secnumdepth}{-\maxdimen} % remove section numbering
\ifLuaTeX
  \usepackage{selnolig}  % disable illegal ligatures
\fi
\IfFileExists{bookmark.sty}{\usepackage{bookmark}}{\usepackage{hyperref}}
\IfFileExists{xurl.sty}{\usepackage{xurl}}{} % add URL line breaks if available
\urlstyle{same}
\hypersetup{
  hidelinks,
  pdfcreator={LaTeX via pandoc}}

\author{}
\date{\vspace{-2.5em}}

\begin{document}

\hypertarget{literature-review-data-technology-investment-and-firm-performance}{%
\section{Literature Review: Data, Technology Investment, and Firm
Performance}\label{literature-review-data-technology-investment-and-firm-performance}}

\textbf{Date:} 2026-03-11 \textbf{Scope:} Effects of customer data
access and data/technology investment on firm revenue, profits, and
productivity --- with emphasis on the airline industry \textbf{Journal
filter:} AER, QJE, JPE, REStud, Econometrica (economics); JF, JFE, RFS
(finance); NBER working paper series \textbf{Note on scope:} The journal
filter is strict. Where a paper appears in a journal outside this list
(e.g., JEL, AEA P\&P, JME, JIE), this is noted explicitly. Papers in the
NBER series that have not yet been published in a target journal are
flagged as working papers.

\begin{center}\rule{0.5\linewidth}{0.5pt}\end{center}

\hypertarget{overview}{%
\subsection{Overview}\label{overview}}

Four questions organize the existing literature relevant to this
project:

\begin{enumerate}
\def\labelenumi{\arabic{enumi}.}
\tightlist
\item
  \textbf{Does data and technology investment raise firm performance?}
  Yes --- in productivity, revenue, and profits --- but the gains are
  conditional on complementary assets and difficult to identify
  causally.
\item
  \textbf{How does customer data specifically enable higher revenue?}
  Through personalization and price discrimination: firms that can infer
  individual willingness to pay charge more to high-type consumers while
  retaining low-type consumers, raising profits.
\item
  \textbf{What does the evidence look like across industries?} The
  airline industry has the most rigorous academic literature on pricing
  mechanisms; retail and e-commerce have the richest data on
  personalization; finance has theoretical models of big data and cost
  of capital; healthcare has causal estimates of IT adoption.
\item
  \textbf{How should one identify these effects?} Causal identification
  is difficult because technology adoption is endogenous. The literature
  uses regulatory shocks (GDPR), field experiments, entry/exit of
  data-intensive competitors, and structural demand models.
\end{enumerate}

\begin{center}\rule{0.5\linewidth}{0.5pt}\end{center}

\hypertarget{section-1-effects-of-data-and-technology-investment-on-firm-performance}{%
\subsection{Section 1: Effects of Data and Technology Investment on Firm
Performance}\label{section-1-effects-of-data-and-technology-investment-on-firm-performance}}

\emph{What do firms gain from data? The evidence on revenue, profits,
and productivity.}

\hypertarget{consensus-findings}{%
\subsubsection{Consensus Findings}\label{consensus-findings}}

\begin{enumerate}
\def\labelenumi{\arabic{enumi}.}
\item
  \textbf{Data and IT investment are associated with higher firm
  productivity, but returns depend on complementarity with existing
  assets.} Firms with pre-existing data assets and data-skilled labor
  earn substantially higher returns from new data technology than firms
  without these complements. There is no universal ``data dividend'' ---
  adoption alone is insufficient.
\item
  \textbf{Personalized pricing, enabled by customer-level data, raises
  firm profits substantially.} The gains from moving from a uniform
  price to an optimized but non-personalized price are large;
  personalization on top of that adds further gains. The welfare
  distribution between firms and consumers is ambiguous and depends on
  the competitive environment.
\item
  \textbf{Privacy regulation --- the loss of data access --- reduces
  firm revenue and productivity, providing a revealed-preference
  estimate of data's value.} GDPR compliance reduced website traffic and
  revenues for data-intensive online firms and measurably reduced data
  storage and computation by EU firms, implying that access to personal
  data was economically meaningful.
\item
  \textbf{Big data disproportionately benefits large firms}, creating a
  feedback loop: larger firms generate more data, which improves their
  forecasts and pricing, which makes them more competitive, which
  generates more data. This mechanism contributes to increasing
  firm-size concentration in data-intensive industries.
\end{enumerate}

\hypertarget{key-papers}{%
\subsubsection{Key Papers}\label{key-papers}}

\begin{itemize}
\item
  \textbf{Dubé \& Misra (2023)} --- \emph{JPE} 131(1): 131--189 ---
  ``Personalized Pricing and Consumer Welfare.'' Randomized controlled
  pricing field experiment on ZipRecruiter. Unexercised market power
  raises profit 55\%; moving to an optimized price adds 19\%; full
  personalization adds 86\% relative to the nonoptimized benchmark. Over
  60\% of consumers benefit from personalization even as total consumer
  surplus falls. The granularity of data and consumer surplus are
  nonmonotonically related. \emph{This is the most directly relevant
  paper to this project's mechanism.}
\item
  \textbf{Bajari, Chernozhukov, Hortaçsu \& Suzuki (2019)} --- \emph{AEA
  Papers \& Proceedings} 109: 33--37 {[}NBER WP 24334{]} --- ``The
  Impact of Big Data on Firm Performance: An Empirical Investigation.''
  Uses proprietary Amazon retail data. Forecast accuracy improves with
  more time-series data (T) but shows diminishing returns; cross-product
  pooling (N) shows near-flat gains. Establishes the empirical framework
  linking data accumulation to operational performance.
\item
  \textbf{Farboodi \& Veldkamp (NBER WP 28427)} --- ``A Model of the
  Data Economy.'' Theoretical model treating data as a production factor
  that helps firms forecast uncertain outcomes and optimize decisions.
  Characterizes data feedback loops, depreciation, and equilibrium data
  prices. Provides the theoretical foundation for thinking about why
  customer data is valuable to a firm.
\item
  \textbf{Goldfarb \& Tucker (2019)} --- \emph{Journal of Economic
  Literature} 57(1): 3--43 {[}not in target journal list, but the
  standard survey{]} --- ``Digital Economics.'' Documents that digital
  technology lowers five key costs: search, replication, transportation,
  tracking, and verification. The reduction in \emph{tracking costs} is
  what enables personalization and data-based price discrimination.
  Canonical reference for contextualizing data investment in the broader
  digital economy.
\end{itemize}

\begin{center}\rule{0.5\linewidth}{0.5pt}\end{center}

\hypertarget{section-2-evidence-by-industry}{%
\subsection{Section 2: Evidence by
Industry}\label{section-2-evidence-by-industry}}

\emph{The same mechanism --- data enables better targeting --- operates
across industries, but the identification strategies and data sources
differ substantially.}

\hypertarget{a.-airlines}{%
\subsubsection{2A. Airlines}\label{a.-airlines}}

\textbf{Consensus:} - Fare variation is large and driven by demand-side
segmentation (business vs.~leisure), not costs. The expected absolute
fare difference between two passengers on the same route is
\textasciitilde36\% of the average fare. - Airlines use screening
devices (advance purchase restrictions, refundability, seat class) to
separate consumer types. Even coarse behavioral signals --- day of
purchase, time-to-departure --- generate measurable fare differences. -
Current pricing captures \textasciitilde77\% of first-best welfare; the
23\% gap arises primarily from \emph{unobservability of passenger type}.
This is the welfare cost of imperfect information about customer
identity --- exactly what customer data acquisition could reduce. -
Dynamic pricing (prices changing over the booking window) is
welfare-superior to uniform pricing in aggregate.

\textbf{Key papers:}

\begin{itemize}
\tightlist
\item
  \textbf{Borenstein \& Rose (1994)} --- \emph{JPE} 102(4): 653--683 ---
  Foundational documentation of price dispersion. Establishes the
  baseline: discrimination, not costs, drives fare variation.
\item
  \textbf{Gerardi \& Shapiro (2009)} --- \emph{JPE} 117(1): 1--37 ---
  Panel FE resolves the Borenstein-Rose puzzle. Competition reduces
  price dispersion within routes; cross-sectional estimates suffer
  omitted variable bias.
\item
  \textbf{Williams (2022)} --- \emph{Econometrica} 90(2): 831--858 ---
  Structural model of dynamic pricing. Dynamic pricing welfare-superior
  to uniform. Decomposes price changes into demand-shock
  vs.~WTP-variation components.
\item
  \textbf{Aryal, Murry \& Williams (2024)} --- \emph{REStud} 91(2):
  641--689 --- Most complete welfare decomposition. 23\% welfare gap
  from private information (unobservable type). Welfare would improve if
  airlines could directly observe passenger type. \emph{Direct
  motivation for the data value channel.}
\end{itemize}

\hypertarget{b.-retail-and-e-commerce}{%
\subsubsection{2B. Retail and
E-Commerce}\label{b.-retail-and-e-commerce}}

\textbf{Consensus:} - Recommendation systems and personalized targeting
raise clicks, conversions, and revenue substantially; banning personal
data from recommendation algorithms sharply reduces sales. - Gains from
data are largest when consumer heterogeneity is high and when the
platform controls the transaction interface (captive channel). - Large
platforms generate more data, which improves their targeting, which
attracts more consumers and generates more data --- creating a
competitive moat.

\textbf{Key papers:}

\begin{itemize}
\tightlist
\item
  \textbf{Bajari, Chernozhukov, Hortaçsu \& Suzuki (2019)} --- \emph{AEA
  P\&P} / NBER WP 24334 --- As above. Amazon retail; data accumulation
  improves forecast accuracy and thereby inventory decisions and
  revenue.
\item
  \textbf{Farboodi, Mihet, Philippon \& Veldkamp (2019)} --- \emph{AEA
  P\&P} 109: 38--42 {[}NBER WP 25515{]} --- ``Big Data and Firm
  Dynamics.'' More data → more data investment → larger firm size
  distribution skewness. Large firms benefit disproportionately from
  data accumulation.
\item
  \textbf{Dubé \& Misra (2023)} --- \emph{JPE} --- As above.
  Retail/online services context (ZipRecruiter).
\end{itemize}

\hypertarget{c.-finance-and-banking}{%
\subsubsection{2C. Finance and Banking}\label{c.-finance-and-banking}}

\textbf{Consensus:} - Big data lowers the cost of capital
\emph{disproportionately for large firms}: larger firms have more
financial history and transaction data, so they benefit more from
data-driven underwriting and valuation. - Fintech lenders using
alternative data (non-FICO signals) extend credit to previously unserved
borrowers and predict default more accurately, suggesting traditional
credit scoring leaves revenue on the table. - The data advantage
compounds: firms with lower cost of capital grow faster, generating more
transactions and more data.

\textbf{Key papers:}

\begin{itemize}
\tightlist
\item
  \textbf{Begenau, Farboodi \& Veldkamp (2018)} --- NBER WP 24550;
  published in \emph{Journal of Monetary Economics} {[}not in target
  journal list{]} --- ``Big Data in Finance and the Growth of Large
  Firms.'' Data lowers cost of capital disproportionately for large
  firms; explains part of the rise in firm-size concentration.
\item
  ⚠️ \emph{Note:} The finance journal literature (JF, JFE, RFS) on data
  value and firm performance is nascent. No confirmed top-3 finance
  paper on data investment → revenue/profits was found in this search.
  This is itself a gap in the literature --- and one reason a paper in
  this space could be publishable in a finance journal.
\end{itemize}

\hypertarget{d.-healthcare}{%
\subsubsection{2D. Healthcare}\label{d.-healthcare}}

\textbf{Consensus:} - Electronic health records (EHR) and health IT
improve care quality and patient outcomes, but productivity effects on
hospitals are mixed and take time to materialize. - The gains from data
in healthcare come from coordination and decision support, not pricing
--- a different channel than in commercial settings.

\textbf{Key papers:}

\begin{itemize}
\tightlist
\item
  \textbf{Lee, McCullough \& Town (2013)} --- NBER WP 18025 --- ``The
  Impact of Health Information Technology on Hospital Productivity.''
  Causal estimates using a value-added production function correcting
  for endogenous input choices. Health IT inputs grew 210\% over the
  study period but contributed only \textasciitilde6\% to value-added
  growth.
\item
  \textbf{Agha (2014)} --- \emph{AEJ: Economic Policy} 6(4) {[}not in
  target journal list{]} --- ``The Effects of Health IT on the Costs and
  Quality of Medical Care.'' HITECH Act adoption as quasi-experiment.
  Mixed productivity effects.
\end{itemize}

\begin{center}\rule{0.5\linewidth}{0.5pt}\end{center}

\hypertarget{section-3-theoretical-foundations-how-does-data-create-value}{%
\subsection{Section 3: Theoretical Foundations --- How Does Data Create
Value?}\label{section-3-theoretical-foundations-how-does-data-create-value}}

\emph{This section is my addition. Understanding the mechanism is
critical for designing empirical tests and interpreting reduced-form
results.}

\hypertarget{consensus-findings-1}{%
\subsubsection{Consensus Findings}\label{consensus-findings-1}}

\begin{enumerate}
\def\labelenumi{\arabic{enumi}.}
\item
  \textbf{Data's value comes from improving predictions, not from data
  itself.} Data helps firms forecast random outcomes (demand, default
  risk, passenger type) more accurately. Better forecasts enable better
  decisions (pricing, inventory, lending). The value is proportional to
  the reduction in forecast error.
\item
  \textbf{Data exhibits non-rivalry and positive network effects, but
  also diminishing returns.} Unlike physical capital, data can be
  replicated at near-zero cost. More data improves forecasts, but with
  diminishing marginal returns (forecast error falls at approximately
  1/√N). This creates a U-shaped dynamic: early data investment has high
  returns; mature firms face diminishing returns but their scale creates
  a moat.
\item
  \textbf{Data is most valuable when consumer heterogeneity is high and
  the firm can act on information asymmetries.} In markets with
  heterogeneous willingness to pay (like airlines), better information
  about individual types enables finer price discrimination. The welfare
  gain from resolving information asymmetry is bounded by the gap
  between actual profits and first-best (Aryal et al.~2024: 23\% of
  first-best welfare).
\end{enumerate}

\hypertarget{key-papers-1}{%
\subsubsection{Key Papers}\label{key-papers-1}}

\begin{itemize}
\tightlist
\item
  \textbf{Farboodi \& Veldkamp (NBER WP 28427)} --- ``A Model of the
  Data Economy.'' Data as a production factor with feedback loops;
  growth model characterizing equilibrium data accumulation.
\item
  \textbf{Dubé \& Misra (2023)} --- \emph{JPE} --- Empirical
  demonstration of the theory: granularity of data determines both the
  magnitude of profit gains and the distribution between firm and
  consumers.
\item
  \textbf{Aryal, Murry \& Williams (2024)} --- \emph{REStud} ---
  Structural quantification of the value of resolving information
  asymmetry. The 23\% welfare gap is the upper bound on the revenue gain
  from perfect customer-type identification.
\end{itemize}

\begin{center}\rule{0.5\linewidth}{0.5pt}\end{center}

\hypertarget{section-4-identification-strategies-and-data-sources}{%
\subsection{Section 4: Identification Strategies and Data
Sources}\label{section-4-identification-strategies-and-data-sources}}

\emph{How has the literature established causal claims? What does each
approach identify?}

\hypertarget{consensus-findings-2}{%
\subsubsection{Consensus Findings}\label{consensus-findings-2}}

\begin{enumerate}
\def\labelenumi{\arabic{enumi}.}
\item
  \textbf{Endogeneity is the central challenge.} Firms that invest more
  in data technology are systematically different (larger,
  faster-growing, more data-intensive). OLS regressions of performance
  on data investment overstate the causal effect. Causal identification
  requires either (a) exogenous variation in data access, (b) field
  experiments, or (c) structural models.
\item
  \textbf{Privacy regulation shocks are the cleanest natural experiments
  for identifying data value.} GDPR (May 2018) and similar regulations
  provide plausibly exogenous variation in data access --- firms serving
  EU customers faced a sudden, externally imposed cost of data
  collection. Studies exploiting this shock find meaningful revenue and
  productivity effects.
\item
  \textbf{Field experiments are feasible in pricing contexts but limited
  to one firm or product.} Dubé \& Misra (2023) is the gold standard:
  randomly assigned prices across customers. Generalizing these
  single-firm results to the industry level requires structural
  assumptions.
\item
  \textbf{Panel FE is the baseline for observational studies, but
  endogeneity of technology adoption remains.} The key lesson from
  airline pricing (Gerardi-Shapiro): cross-sectional estimates suffer
  omitted variable bias; within-unit (route, firm) variation is
  necessary. Instrument for technology investment timing if possible.
\end{enumerate}

\hypertarget{identification-approaches}{%
\subsubsection{Identification
Approaches}\label{identification-approaches}}

\begin{longtable}[]{@{}
  >{\raggedright\arraybackslash}p{(\columnwidth - 6\tabcolsep) * \real{0.1667}}
  >{\raggedright\arraybackslash}p{(\columnwidth - 6\tabcolsep) * \real{0.3167}}
  >{\raggedright\arraybackslash}p{(\columnwidth - 6\tabcolsep) * \real{0.2667}}
  >{\raggedright\arraybackslash}p{(\columnwidth - 6\tabcolsep) * \real{0.2500}}@{}}
\toprule\noalign{}
\begin{minipage}[b]{\linewidth}\raggedright
Strategy
\end{minipage} & \begin{minipage}[b]{\linewidth}\raggedright
What it identifies
\end{minipage} & \begin{minipage}[b]{\linewidth}\raggedright
Key assumption
\end{minipage} & \begin{minipage}[b]{\linewidth}\raggedright
Best examples
\end{minipage} \\
\midrule\noalign{}
\endhead
\bottomrule\noalign{}
\endlastfoot
Randomized experiment & ATE of personalization/data on profits & Random
assignment of prices/treatment & Dubé \& Misra (2023) \\
Regulatory shock (GDPR, CCPA) & Effect of losing data access on revenue
& Parallel trends; regulation is exogenous to firm-level outcomes &
Aridor, Che \& Salz (NBER); Goldfarb-Tucker \\
Entry/exit of data-intensive competitor & Competitive effect of data
advantage & Entry is orthogonal to route-level demand trends & Gap for
this project \\
Technology adoption timing (DiD/event study) & Effect of investment on
pricing outcomes & Adoption timing uncorrelated with demand shocks &
\textbf{This paper's proposed design} \\
Structural demand model & Welfare counterfactuals; markup decomposition
& Functional form and equilibrium assumptions & Williams (2022); Aryal
et al.~(2024) \\
\end{longtable}

\hypertarget{data-sources-relevant-to-this-project}{%
\subsubsection{Data Sources Relevant to This
Project}\label{data-sources-relevant-to-this-project}}

\begin{longtable}[]{@{}
  >{\raggedright\arraybackslash}p{(\columnwidth - 6\tabcolsep) * \real{0.1538}}
  >{\raggedright\arraybackslash}p{(\columnwidth - 6\tabcolsep) * \real{0.2885}}
  >{\raggedright\arraybackslash}p{(\columnwidth - 6\tabcolsep) * \real{0.2500}}
  >{\raggedright\arraybackslash}p{(\columnwidth - 6\tabcolsep) * \real{0.3077}}@{}}
\toprule\noalign{}
\begin{minipage}[b]{\linewidth}\raggedright
Source
\end{minipage} & \begin{minipage}[b]{\linewidth}\raggedright
What it covers
\end{minipage} & \begin{minipage}[b]{\linewidth}\raggedright
Key strength
\end{minipage} & \begin{minipage}[b]{\linewidth}\raggedright
Key limitation
\end{minipage} \\
\midrule\noalign{}
\endhead
\bottomrule\noalign{}
\endlastfoot
BTS DB1B & U.S. domestic itinerary-level fares, quarterly, 10\% sample &
Large; standard in literature; covers all carriers & No purchase date;
quarterly frequency; transacted (not offered) fares \\
BTS T-100 & U.S. route-level traffic and revenue, monthly & Monthly
frequency; complete census & Aggregate; no fare distribution \\
ATPCO & Filed fares by fare class, real-time & Offered prices; fare
rules (advance purchase, refundability) & Requires paid access; very
high volume \\
Scraped airline websites & Offered prices over booking window &
Booking-window dynamics & Snapshot; route/airline selection; not
transacted \\
Airline 10-K / SEC filings & IT capex, technology investment, loyalty
program stats & Causal variation candidate; public & Qualitative;
inconsistent disclosure across carriers and years \\
Loyalty program investor materials & Enrollment rates, miles sold,
credit card partnerships & Data asset proxy; publicly disclosed by
Delta, United, American & Carrier-level; not route-level \\
\end{longtable}

\begin{center}\rule{0.5\linewidth}{0.5pt}\end{center}

\hypertarget{section-5-open-questions-and-this-papers-contribution}{%
\subsection{Section 5: Open Questions and This Paper's
Contribution}\label{section-5-open-questions-and-this-papers-contribution}}

\begin{enumerate}
\def\labelenumi{\arabic{enumi}.}
\item
  \textbf{The causal effect of data investment on pricing outcomes is
  unidentified.} All existing causal work in the airline space uses
  \emph{competition} as the source of variation. No paper uses variation
  in \emph{data technology investment timing} as the treatment. This is
  the gap.
\item
  \textbf{Customer data acquisition and technology investment are
  distinct, unstudied channels.} Technology improves the optimization
  algorithm; customer data improves the inputs. A loyalty program
  expansion increases the data asset; a revenue management system
  upgrade improves how that asset is used. The academic literature has
  not separated these.
\item
  \textbf{No top-finance-journal paper exists on data value in the
  airline industry.} The finance literature has models of big data and
  cost of capital but nothing on airlines, pricing, or the revenue
  mechanism. A paper that connects causal pricing evidence to a firm
  value / return on data investment framing would be well-positioned for
  JF/JFE/RFS.
\item
  \textbf{Welfare vs.~revenue.} The IO literature focuses on welfare.
  Finance and corporate strategy care about \emph{revenue and profit}.
  Our paper operates in revenue space --- more tractable with
  reduced-form methods and directly relevant to valuation.
\item
  \textbf{Heterogeneous returns.} Tambe (2014) finds data returns are
  highest in data-intensive industries. The analog: are data technology
  returns larger on routes with higher demand heterogeneity (more to
  gain from finer segmentation)? This cross-sectional heterogeneity
  would be a testable and novel prediction.
\end{enumerate}

\begin{center}\rule{0.5\linewidth}{0.5pt}\end{center}

\hypertarget{bibtex-entries}{%
\subsection{BibTeX Entries}\label{bibtex-entries}}

\begin{Shaded}
\begin{Highlighting}[]
\VariableTok{@article}\NormalTok{\{}\OtherTok{DubeMisra2023}\NormalTok{,}
  \DataTypeTok{author}\NormalTok{  = \{Dub}\CharTok{\textbackslash{}\textquotesingle{}}\NormalTok{\{e\}, Jean{-}Pierre and Misra, Sanjog\},}
  \DataTypeTok{title}\NormalTok{   = \{Personalized Pricing and Consumer Welfare\},}
  \DataTypeTok{journal}\NormalTok{ = \{Journal of Political Economy\},}
  \DataTypeTok{volume}\NormalTok{  = \{131\},}
  \DataTypeTok{number}\NormalTok{  = \{1\},}
  \DataTypeTok{pages}\NormalTok{   = \{131{-}{-}189\},}
  \DataTypeTok{year}\NormalTok{    = \{2023\}}
\NormalTok{\}}

\VariableTok{@article}\NormalTok{\{}\OtherTok{BorensteinRose1994}\NormalTok{,}
  \DataTypeTok{author}\NormalTok{  = \{Borenstein, Severin and Rose, Nancy L.\},}
  \DataTypeTok{title}\NormalTok{   = \{Competition and Price Dispersion in the \{U.S.\} Airline Industry\},}
  \DataTypeTok{journal}\NormalTok{ = \{Journal of Political Economy\},}
  \DataTypeTok{volume}\NormalTok{  = \{102\},}
  \DataTypeTok{number}\NormalTok{  = \{4\},}
  \DataTypeTok{pages}\NormalTok{   = \{653{-}{-}683\},}
  \DataTypeTok{year}\NormalTok{    = \{1994\}}
\NormalTok{\}}

\VariableTok{@article}\NormalTok{\{}\OtherTok{GerardiShapiro2009}\NormalTok{,}
  \DataTypeTok{author}\NormalTok{  = \{Gerardi, Kristopher and Shapiro, Adam Hale\},}
  \DataTypeTok{title}\NormalTok{   = \{Does Competition Reduce Price Dispersion? \{New\} Evidence from the Airline Industry\},}
  \DataTypeTok{journal}\NormalTok{ = \{Journal of Political Economy\},}
  \DataTypeTok{volume}\NormalTok{  = \{117\},}
  \DataTypeTok{number}\NormalTok{  = \{1\},}
  \DataTypeTok{pages}\NormalTok{   = \{1{-}{-}37\},}
  \DataTypeTok{year}\NormalTok{    = \{2009\}}
\NormalTok{\}}

\VariableTok{@article}\NormalTok{\{}\OtherTok{Williams2022}\NormalTok{,}
  \DataTypeTok{author}\NormalTok{  = \{Williams, Kevin R.\},}
  \DataTypeTok{title}\NormalTok{   = \{The Welfare Effects of Dynamic Pricing: \{Evidence\} from Airline Markets\},}
  \DataTypeTok{journal}\NormalTok{ = \{Econometrica\},}
  \DataTypeTok{volume}\NormalTok{  = \{90\},}
  \DataTypeTok{number}\NormalTok{  = \{2\},}
  \DataTypeTok{pages}\NormalTok{   = \{831{-}{-}858\},}
  \DataTypeTok{year}\NormalTok{    = \{2022\}}
\NormalTok{\}}

\VariableTok{@article}\NormalTok{\{}\OtherTok{AryalMurryWilliams2024}\NormalTok{,}
  \DataTypeTok{author}\NormalTok{  = \{Aryal, Gaurab and Murry, Charles and Williams, Jonathan W.\},}
  \DataTypeTok{title}\NormalTok{   = \{Price Discrimination in International Airline Markets\},}
  \DataTypeTok{journal}\NormalTok{ = \{Review of Economic Studies\},}
  \DataTypeTok{volume}\NormalTok{  = \{91\},}
  \DataTypeTok{number}\NormalTok{  = \{2\},}
  \DataTypeTok{pages}\NormalTok{   = \{641{-}{-}689\},}
  \DataTypeTok{year}\NormalTok{    = \{2024\}}
\NormalTok{\}}

\VariableTok{@article}\NormalTok{\{}\OtherTok{BajariChernozhukovHortacsuSuzuki2019}\NormalTok{,}
  \DataTypeTok{author}\NormalTok{  = \{Bajari, Patrick and Chernozhukov, Victor and Horta}\CharTok{\textbackslash{}c}\NormalTok{\{c\}su, Ali and Suzuki, Junichi\},}
  \DataTypeTok{title}\NormalTok{   = \{The Impact of Big Data on Firm Performance: \{An\} Empirical Investigation\},}
  \DataTypeTok{journal}\NormalTok{ = \{AEA Papers and Proceedings\},}
  \DataTypeTok{volume}\NormalTok{  = \{109\},}
  \DataTypeTok{pages}\NormalTok{   = \{33{-}{-}37\},}
  \DataTypeTok{year}\NormalTok{    = \{2019\},}
  \DataTypeTok{note}\NormalTok{    = \{Also NBER Working Paper No.\textasciitilde{}24334\}}
\NormalTok{\}}

\VariableTok{@article}\NormalTok{\{}\OtherTok{FarboodiMihetPhilipponVeldkamp2019}\NormalTok{,}
  \DataTypeTok{author}\NormalTok{  = \{Farboodi, Maryam and Mihet, Roxana and Philippon, Thomas and Veldkamp, Laura\},}
  \DataTypeTok{title}\NormalTok{   = \{Big Data and Firm Dynamics\},}
  \DataTypeTok{journal}\NormalTok{ = \{AEA Papers and Proceedings\},}
  \DataTypeTok{volume}\NormalTok{  = \{109\},}
  \DataTypeTok{pages}\NormalTok{   = \{38{-}{-}42\},}
  \DataTypeTok{year}\NormalTok{    = \{2019\},}
  \DataTypeTok{note}\NormalTok{    = \{Also NBER Working Paper No.\textasciitilde{}25515\}}
\NormalTok{\}}

\VariableTok{@techreport}\NormalTok{\{}\OtherTok{FarboodiVeldkamp2021}\NormalTok{,}
  \DataTypeTok{author}\NormalTok{      = \{Farboodi, Maryam and Veldkamp, Laura\},}
  \DataTypeTok{title}\NormalTok{       = \{A Model of the Data Economy\},}
  \DataTypeTok{institution}\NormalTok{ = \{National Bureau of Economic Research\},}
  \DataTypeTok{number}\NormalTok{      = \{28427\},}
  \DataTypeTok{year}\NormalTok{        = \{2021\},}
  \DataTypeTok{type}\NormalTok{        = \{Working Paper\}}
\NormalTok{\}}

\VariableTok{@techreport}\NormalTok{\{}\OtherTok{BegenauFarboodiVeldkamp2018}\NormalTok{,}
  \DataTypeTok{author}\NormalTok{      = \{Begenau, Juliane and Farboodi, Maryam and Veldkamp, Laura\},}
  \DataTypeTok{title}\NormalTok{       = \{Big Data in Finance and the Growth of Large Firms\},}
  \DataTypeTok{institution}\NormalTok{ = \{National Bureau of Economic Research\},}
  \DataTypeTok{number}\NormalTok{      = \{24550\},}
  \DataTypeTok{year}\NormalTok{        = \{2018\},}
  \DataTypeTok{type}\NormalTok{        = \{Working Paper\},}
  \DataTypeTok{note}\NormalTok{        = \{Published in Journal of Monetary Economics\}}
\NormalTok{\}}

\VariableTok{@article}\NormalTok{\{}\OtherTok{GoldfarbTucker2019}\NormalTok{,}
  \DataTypeTok{author}\NormalTok{  = \{Goldfarb, Avi and Tucker, Catherine\},}
  \DataTypeTok{title}\NormalTok{   = \{Digital Economics\},}
  \DataTypeTok{journal}\NormalTok{ = \{Journal of Economic Literature\},}
  \DataTypeTok{volume}\NormalTok{  = \{57\},}
  \DataTypeTok{number}\NormalTok{  = \{1\},}
  \DataTypeTok{pages}\NormalTok{   = \{3{-}{-}43\},}
  \DataTypeTok{year}\NormalTok{    = \{2019\},}
  \DataTypeTok{note}\NormalTok{    = \{Not in target journal list; included as standard survey reference\}}
\NormalTok{\}}

\VariableTok{@techreport}\NormalTok{\{}\OtherTok{LeeMcCulloughTown2013}\NormalTok{,}
  \DataTypeTok{author}\NormalTok{      = \{Lee, Jinhyung and McCullough, Jeffrey S. and Town, Robert J.\},}
  \DataTypeTok{institution}\NormalTok{ = \{National Bureau of Economic Research\},}
  \DataTypeTok{number}\NormalTok{      = \{18025\},}
  \DataTypeTok{title}\NormalTok{       = \{The Impact of Health Information Technology on Hospital Productivity\},}
  \DataTypeTok{type}\NormalTok{        = \{Working Paper\},}
  \DataTypeTok{year}\NormalTok{        = \{2012\}}
\NormalTok{\}}
\end{Highlighting}
\end{Shaded}

\begin{center}\rule{0.5\linewidth}{0.5pt}\end{center}

\hypertarget{verification-notes}{%
\subsection{Verification Notes}\label{verification-notes}}

\begin{itemize}
\tightlist
\item
  \textbf{Dubé \& Misra (2023):} Confirmed JPE 131(1) --- verify exact
  pages.
\item
  \textbf{Aryal et al.~(2024):} Confirmed REStud 91(2) --- verify exact
  pages.
\item
  \textbf{Williams (2022):} Confirmed Econometrica 90(2):831--858.
\item
  \textbf{Bajari et al.~(2019):} Confirmed AEA P\&P 109:33--37. Note
  this is \emph{not} the main AER --- it is the conference proceedings
  volume, which is published alongside the AER but is lighter in review
  standards.
\item
  \textbf{Farboodi \& Veldkamp NBER 28427:} Working paper as of last
  search; confirm if published in a target journal.
\item
  \textbf{Begenau, Farboodi, Veldkamp:} Published in Journal of Monetary
  Economics --- \emph{not} JF/JFE/RFS. Flagged accordingly.
\item
  \textbf{Goldfarb \& Tucker (2019):} JEL --- not on the target list,
  but the standard survey for digital economics.
\item
  \textbf{No top-3 finance journal (JF/JFE/RFS) papers} on data
  investment and airline pricing or personalized pricing were
  identified. This gap is informative for positioning the paper.
\end{itemize}

\end{document}
